\scanpage{23}
\mspage{21}
% archive sheet number: 22
% hand: Grothendieck
% sheet_kind: cancelled_draft
% element: text | top, lines 1–4 | on constate immédiatement
on constate immédiatement qu'\deleted{\unsure{il}} il \unsure{se trouve} $H(\overline R)$ d'indice fini dans $G^\circ(\overline R)$, ce qui implique que $H=G^\circ$; \deleted{donc} comme $H$ est unipotent, il en est de même de $G^\circ$.
% element: margin note | upper left, slanted | dégager un lemme
% slanted
\margin{\unsure{dégager} un lemme, avec le corollaire $E^U=E^{G^\circ}$.}

% element: text | upper-middle, lines 5–7 | 5.4 Supposons maintenant
5.4. Supposons maintenant que $\pi$ opère de façon quasi-unipotente, (\deleted{donc}\inserted{i.e.} $G$ est unipotent) et \uline{supposons que $R$ soit de car. nulle}.

% element: text | cancelled block, lines 7–16 | Si U est un sous-groupe
\deleted{Si $U$ est un sous-groupe \deleted{\unsure{ouvert}}\inserted{invariant} de $\pi$ qui opère de façon unipotente (p.ex $u^{-1}(G^\circ(R))$), on peut former \deleted{\unsure{$F$}} $E^U$, sur lequel $\pi/U$ opère. Je dis que $E^U$ ne dépend pas du choix de $U$; en effet, il suffit de voir que $E^U=E^{U'}$ si $U'\subset U$, $U'$ \uline{invariant}. Or $\pi/U'$ opère sur $E^{U'}$, donc $U/U'$ y opère; c'est un groupe fini qui opère de façon unipotente, donc (puisqu'on est en car nulle) il opère \unsure{trivialt|trivialmt}. De plus, $E^U$ \deleted{\unsure{soit}} \inserted{inv. dans $\pi$ (car on peut supposer $U$ invariant dans $\pi$).}}
% element: text | cancelled line, below line 16 | Ceci nous montre
\deleted{\deleted{Ceci nous montre en particulier que}}
% element: display | cancelled block, line 17 | E^U = E^G^0
\[
\deleted{E^U=E^{G^\circ}.}
\]
% element: text | cancelled block, lines 17–18 | Nous pouvons définir
\deleted{Nous pouvons définir une filtration croissante de $E$, en posant}
% element: display | cancelled block, lines 19–20 | E^(0) = 0
\[
\begin{aligned}
&\deleted{E_{\unsure{0|a}}^{(0)}=0,\qquad E^{(i+1)}=\text{image inverse de}}\\
&\deleted{\deleted{\unsure{(E/E^i)^U=(E/E^i)^{G^\circ}}}}\\
&\deleted{\inserted{(E/E^i)^U=(E/E^i)^{G^\circ}\quad\text{dans }E}.}
\end{aligned}
\]
% element: text | cancelled bottom block, lines 21–24 | L'hypothèse
\deleted{L'hypothèse que \deleted{\unsure{$U$ opère}} $G^\circ$ opère de façon unipotente implique que cette filtration est exhaustive. \deleted{\illegible{1}} D'autre part, $U$ opère \deleted{\unsure{de façon unipotente}}\inserted{trivialmt} sur le gradué associé, par construction. On}
